\chapter{Architecture de communication}

Ce chapitre analyse les protocoles et les flux de donn\'ees du syst\`eme Andon, en justifiant les choix de communication et en d\'ecrivant les m\'ecanismes d'interconnexion entre les diff\'erents composants.

\section{Protocoles de communication utilis\'es}

\subsection{MQTT : protocole principal}

MQTT (Message Queuing Telemetry Transport) est le protocole principal du syst\`eme, assurant la communication entre les capteurs ESP32 et le serveur backend \cite{banks2014mqtt}.

\subsubsection{Justification du choix}

Le choix de MQTT est justifi\'e par :

\begin{itemize}
    \item L\'eg\`eret\'e des messages : en-t\^etes minimaux (2 octets), adapt\'e aux microcontr\^oleurs.
    \item Le mod\`ele publish/subscribe : d\'e-couplage entre producteurs et consommateurs de donn\'ees.
    \item Les niveaux de QoS : garantie de livraison configurable selon la criticit\'e.
    \item La reconnexion automatique : le client se reconnecte apr\`es une d\'econnexion.
    \item La large adoption : support\'e par toutes les plateformes IoT.
\end{itemize}

\subsubsection{Configuration du broker}

Le broker MQTT Mosquitto est configur\'e pour le syst\`eme Andon :

\begin{lstlisting}[caption={Configuration du broker Mosquitto},label={lst:mosquitto_config}]
listener 1883 localhost
listener 8883
certfile /etc/mosquitto/certs/server.crt
keyfile /etc/mosquitto/certs/server.key
allow_anonymous false
password_file /etc/mosquitto/passwd
max_connections -1
max_inflight_messages 20
max_queued_messages 1000
\end{lstlisting}

\subsection{HTTP/HTTPS : API REST}

Le protocole HTTP est utilis\'e pour les communications entre le frontend et le backend via l'API REST :

\begin{itemize}
    \item Les requ\^etes GET pour la lecture de donn\'ees.
    \item Les requ\^etes POST pour la cr\'eation de ressources.
    \item Les requ\^etes PUT pour la mise \`a jour de ressources.
    \item Les requ\^etes DELETE pour la suppression de ressources.
\end{itemize}

En environnement de production, le HTTPS est obligatoire pour chiffrer les communications.

\subsection{WebSocket (Socket.IO)}

Socket.IO est utilis\'e pour les communications temps r\'eel bidirectionnelles entre le serveur et le frontend :

\begin{itemize}
    \item Diffusion instantan\'ee des \'ev\'enements Andon.
    \item Mise \`a jour de l'\`etat des lignes en direct.
    \item Notifications visuelles en temps r\'eel.
    \item Fallback automatique sur HTTP long-polling si WebSocket n'est pas disponible.
\end{itemize}

\section{Structure des topics MQTT}

\subsection{Hierarchie des topics}

La hi\'erarchie des topics MQTT refl\`ete l'organisation physique de l'atelier :

\begin{lstlisting}[caption={Hi\'erarchie des topics MQTT},label={lst:topics_hierarchie}]
# Alertes Andon (du capteur au serveur)
sews/andon/1/alert
sews/andon/2/alert
...
sews/andon/8/alert

# Status des capteurs
sews/capteurs/1/status
sews/capteurs/2/status
...
sews/capteurs/8/status

# Commandes retour (du serveur au capteur)
sews/command/1/led
sews/command/2/led
...
sews/command/8/led
\end{lstlisting}

\subsection{Format des messages}

Les messages MQTT utilisent le format JSON pour faciliter le traitement :

\begin{lstlisting}[caption={Format du message d'alerte Andon},label={lst:message_format}]
{
  "ligneId": 1,
  "timestamp": "2026-04-15T10:30:00.000Z",
  "type": "panne_electrique",
  "urgence": "haute",
  "operateur": "op_t1_01",
  "description": "Court-circuit detecte 
    sur testeur 3"
}
\end{lstlisting}

\section{Flux de donn\'ees d\'etaill\'e}

\subsection{Flux d'un \'ev\'enement Andon}

Le parcours complet d'un \'ev\'enement Andon suit les \'etapes suivantes :

\begin{enumerate}
    \item \textbf{D\'etection :} l'op\'erateur appuie sur le bouton physique.
    \item \textbf{Traitement local :} l'ESP32 d\'etecte l'appui, applique un d\'ebounce de 200 ms et g\'en\`ere un message JSON.
    \item \textbf{Publication MQTT :} le message est publi\'e sur le topic \texttt{sews/andon/\{ligne\_id\}/alert} avec QoS 1.
    \item \textbf{Routage :} le broker MQTT redistribue le message aux abonn\'es (serveur backend).
    \item \textbf{Traitement serveur :} le serveur Node.js re\c{c}oit le message, valide les donn\'ees et cr\'e un enregistrement dans PostgreSQL.
    \item \textbf{Diffusion temps r\'eel :} le serveur diffuse l'\'ev\'enement via Socket.IO aux clients connect\'es.
    \item \textbf{Affichage :} le frontend Angular re\c{c}oit la notification et met \`a jour le tableau de bord.
    \item \textbf{Notification :} un e-mail est envoy\'e au technicien de garde.
    \item \textbf{Indicateur visuel :} le serveur publie une commande sur \texttt{sews/command/\{ligne\_id\}/led} pour allumer la LED rouge.
\end{enumerate}

\subsection{Flux de r\'esolution}

Lorsqu'un technicien r\'e\'esout une panne :

\begin{enumerate}
    \item Le technicien acc\`ede \`a l'\'ev\'enement dans le frontend.
    \item Il saisit le compte-rendu d'intervention et valide la r\'e-solution.
    \item Le frontend envoie une requ\^ete PUT \`a l'API REST (\texttt{PUT /api/evenements/:id}).
    \item Le serveur met \`a jour l'enregistrement dans PostgreSQL (statut = resolu, date\_fin = now).
    \item Le serveur publie la commande de r\'e-solution sur MQTT.
    \item L'ESP32 re\c{c}oit la commande et\'eint la LED rouge, allume la LED verte.
    \item Le frontend affiche la ligne comme r\'e-solue.
    \item Une notification de r\'e-solution est envoy\'ee.
\end{enumerate}

\section{Gestion de la s\'ecurit\'e des communications}

\subsection{S\'ecurit\'e MQTT}

La s\'ecurit\'e de la communication MQTT est assur\'ee par :

\begin{itemize}
    \item \textbf{Authentification :} chaque client MQTT doit fournir un identifiant et un mot de passe.
    \item \textbf{TLS :} chiffrement des communications sur le port 8883.
    \item \textbf{ACL :} contr\^ole d'acc\`es par topic (chaque n\oe{}ud ne publie que sur son topic).
    \item \textbf{Persistance :} les messages QoS 1 sont conserv\'es jusqu'\`a r\'e-ception de l'accus\'e.
\end{itemize}

\subsection{S\'ecurit\'e HTTP}

La s\'ecurit\'e de l'API REST est assur\'ee par :

\begin{itemize}
    \item \textbf{HTTPS :} chiffrement TLS des communications.
    \item \textbf{JWT :} tokens d'authentification sign\'es avec expiration.
    \item \textbf{Validation :} v\'erification des donn\'ees d'entr\'ee.
    \item \textbf{Rate limiting :} limitation du nombre de requ\^etes.
\end{itemize}

\section{Gestion des erreurs et tol\'erance aux pannes}

\subsection{Strat\'egie de reconnexion}

Le client MQTT int\`egre un m\'ecanisme de reconnexion automatique :

\begin{itemize}
    \item Tentative de reconnexion toutes les 5 secondes en cas de d\'econnexion.
    \item Backoff exponentiel avec un d\'elai maximum de 60 secondes.
    \item Conservation des messages en attente pendant la d\'econnexion.
    \item Re-souscription aux topics lors de la reconnexion.
\end{itemize}

\subsection{Persistence des donn\'ees}

En cas de d\'efaut du serveur :

\begin{itemize}
    \item Les messages MQTT non trait\'es sont conserv\'es dans le broker (QoS 1).
    \item PostgreSQL assure la durabilit\'e des donn\'ees enregistr\'ees (WAL).
    \item Le frontend affiche les donn\'ees en cache en attendant la reconnexion.
\end{itemize}

\section{Synth\`ese du chapitre}

L'architecture de communication du syst\`eme repose sur trois protocoles compl\'ementaires : MQTT pour la couche capteur, HTTP/REST pour l'API, et Socket.IO pour le temps r\'eel. La hi\'erarchie des topics MQTT, les flux de donn\'ees d\'etaill\'es et les m\'ecanismes de s\'ecurit\'e forment un syst\`eme coh\'erent et fiable, capable de g\'erer les contraintes de l'environnement industriel.
